
\section{An optimal universal radial carrier by source-dependent seven-ray splitting}
\label{sec:xxviii-sharp-source-dependent}
The linear-programming optimum of Edition XXVII concerns a
source-independent \emph{fixed} coefficient splitting. Allowing the
isotropic gauge and the redundant seven-ray coefficient to depend
smoothly on the source tensor yields a smaller universal radial-carrier
cost. This result is purely a symmetric-tensor pressure-quotient
lemma on a periodic-axial annulus, not a source-matched unforced
Navier--Stokes theorem.

\begin{theorem}[Sharp universal annular carrier and smooth seven-ray selection]
\label{thm:xxviii-sharp-universal}
Let $Q(r,z)$ be a smooth symmetric co-rotating $3\times3$ tensor
with compact radial support away from the symmetry axis on
$\mathbb R^2\times\mathbb S^1_z$, or in an embedded transverse
disk on $\mathbb T^3$. Let
\[
M=\max_{i,j}\|Q_{ij}\|_\infty,\qquad \mu>0,
\]
and take a smooth radial cutoff $0\le\chi\le1$ equal to one
on a neighborhood of the full radial support of $Q$.
There exist \emph{explicit smooth} coefficients $d_1,\ldots,d_7$
and a smooth scalar $a=a(r,z)$, all radially compactly
supported, such that
\begin{equation}\label{eq:xxviii-exact-stress}
\boxed{
d_j\ge\mu\chi,\qquad
S:=\sum_{j=1}^7d_jv_j\otimes v_j
=Q+aI+B\chi e_r\otimes e_r,\quad
B=\frac{85}{8}\mu+\frac{343}{4}M.
}
\end{equation}
The rays have aperture $1/4$, exactly as in Editions XXV--XXVII.
The reconstructed stress satisfies the explicit bound
\begin{equation}\label{eq:xxviii-supbound}
\boxed{\quad
\|S\|_{L^\infty,\mathrm{entry}}\le90M+11\mu.\quad}
\end{equation}
Moreover,
\begin{equation}\label{eq:xxviii-projected}
\nabla\cdot(S-Q)
=\nabla\left[a+B\chi+B\int_0^r\frac{\chi(s)}s\,ds\right],
\qquad
\mathbb P\nabla\cdot S=\mathbb P\nabla\cdot Q.
\end{equation}
The carrier coefficient $343/4$ is \emph{sharp} for the universal
independent-entry box $|Q_{ij}|\le M$, when the radial coefficient
$B$ is uniform over the source family and the isotropic scalar
$a(Q)$ and all seven ray coefficients may depend on $Q$.
It is not claimed sharp for the smaller class of positive
semidefinite Reynolds covariance matrices or for a fixed
physical A.11/A.19 source.
\end{theorem}

\begin{proof}[Constructive exact-rational proof]
Represent the independent symmetric entries of $Q$ by
$q=(Q_{rr},Q_{\theta\theta},Q_{zz},
Q_{r\theta},Q_{rz},Q_{\theta z})$.
Write $s=(s_1,\ldots,s_6)\in\{-1,+1\}^6$, with the
$64$ vertices ordered lexicographically.
For each $s$ the accompanying frozen table
\texttt{XXVIII\_64\_CORNER\_EXACT\_CERTIFICATES.txt}
gives rational numbers
$(a_s,d_{s,1},\ldots,d_{s,7})$ with
$d_{s,j}\ge0$, $a_s\le13/4$, and the six \emph{exact} identities
\begin{equation}\label{eq:xxviii-vertices}
\sum_{j=1}^7d_{s,j}v_j\otimes v_j
=Q_s+a_sI+\frac{343}{4}e_r\otimes e_r,
\end{equation}
where $Q_s$ is the symmetric matrix with independent
entries $s_i$. These are not floating-point certificates:
all $64$ rows are independently checked with rational
arithmetic by \texttt{verify\_xxviii\_sharp\_cone.py}.

For $M>0$ put $x_i=q_i/M\in[-1,1]$ and use the nonnegative
multilinear barycentric polynomials
\[
w_s(x)=2^{-6}\prod_{i=1}^6(1+s_i x_i),\qquad
\sum_s w_s=1,\qquad
\sum_s w_s s_i=x_i.
\]
Let
$d^0=(1,1,3,2,2,1,1)$ be the source-free ray
coefficient vector. It obeys
\[
\sum_j d_j^0v_j\otimes v_j
=\frac38 I+\frac{85}{8}e_r\otimes e_r.
\]
Define
\begin{align*}
a(r,z)&=\chi(r)\left[\frac38\mu+
 M\sum_s w_s(q/M)a_s\right],\\
d_j(r,z)&=\chi(r)\left[\mu d_j^0+
 M\sum_s w_s(q/M)d_{s,j}\right].
\end{align*}
All functions are smooth, because each $w_s$ is a polynomial
and $Q$ is smooth, and the radial cutoff is identically one
where $Q\ne0$. Weighted addition of
\eqref{eq:xxviii-vertices} proves
\eqref{eq:xxviii-exact-stress} everywhere, including the
radial collar where $Q=0$ and $0<\chi<1$.
The nonnegative weights and certificates prove $d_j\ge\mu\chi$.
For $M=0$, use only the source-free formula. Since
$0\le a_s\le13/4$, one has
$0\le a\le(3\mu/8+13M/4)\chi$.
Combining with $B$ proves
\eqref{eq:xxviii-supbound}; all off-diagonal entries
of $S$ equal those of $Q$. The pressure-equivalence
formula follows by direct tensor differentiation of the
radial-only carrier and the isotropic tensor.

For \emph{sharpness}, define a dual linear functional on
symmetric matrices by
\[
L(H)=-H_{rr}-15H_{\theta\theta}+16H_{zz}
-\frac{31}{4}H_{r\theta}-46H_{\theta z}.
\]
Then $L(I)=0$, $L(e_r\otimes e_r)=-1$, and the exact
seven ray values are
\[
L(v_j\otimes v_j)
=\left(-1,-\frac{31}{8},0,0,0,-\frac{23}{4},0\right).
\]
Hence every strict ray tensor with $d_j\ge\mu$
satisfies $L(S)\le-85\mu/8$. If
$S=Q+aI+B e_r\otimes e_r$, it follows that
\[
B\ge L(Q)+\frac{85}{8}\mu.
\]
Take the admissible signed box corner
$q=(-M,-M,M,-M,-M,-M)$: then
$L(Q)=343M/4$. This proves that the displayed
$B$ is the smallest possible universal radial coefficient,
with a matching feasible exact rational corner certificate.
\end{proof}

\paragraph{Interpretation and remaining source theorem.}
The improvement is from the Edition XXVII $471M/4$
carrier slope to $343M/4$---a $27.2\%$ reduction---and
from $122M+11\mu$ to $90M+11\mu$ in the displayed
entrywise stress bound. Unlike a single local LP solution,
the multilinear formula constructs smooth globally
cutoff-compatible coefficients. However, the original
OpenAI source's actual A.11/A.19 profiles, physical forcing,
three-dimensional compact shell virtual-work conditions,
space--time stress budget, nonlinear all-harmonic correction
and zero-force singularity theorem remain \emph{unproved}.
A finite Lean scalar dual is supplied; the 64-corner
interpolation and physical PDE are not yet formalized in
Lean. The fully unforced Navier--Stokes problem remains open.
